\section*{Chapter \ref{ch:b3:planetary-atmospheres} --- Planetary Atmospheres and Habitability}

\begin{solution}{exo:b3:planetary-atmospheres:1}
$H = RT/Mg$: Venus $8.314\times737/(0.04401\times8.87) = \qty{15.7}{km}$;
the Earth $8.314\times288/(0.02897\times9.80) = \qty{8.4}{km}$; Mars
$8.314\times214/(0.04334\times3.71) = \qty{11.1}{km}$; Titan
$8.314\times94/(0.0280\times1.35) = \qty{20.7}{km}$. Titan's atmosphere,
cold as it is, is the most extended: its gravity is seven times weaker
than the Earth's.
\end{solution}

\begin{solution}{exo:b3:planetary-atmospheres:2}
$\lambda = GMm/kTr = \num{6.674e-11}\times\num{5.97e24}\times\num{1.674e-27}
/(\num{1.381e-23}\,T\times\num{6.878e6})$: 8.78 at \qty{800}{K}, 4.39 at
\qty{1600}{K}. The escaping fractions are $9.78\,e^{-8.78} =
\num{1.5e-3}$ and $5.39\,e^{-4.39} = \num{6.7e-2}$: the solar cycle
changes the escaping fraction of hydrogen by a factor of about 44.
\end{solution}

\begin{solution}{exo:b3:planetary-atmospheres:3}
(a)~$d = S_\mathrm{eff}^{-1/2}$: \omterm{def:b3:planetary-atmospheres:runaway}{moist greenhouse} \qty{0.993}{au}, runaway
\qty{0.975}{au}, maximum greenhouse \qty{1.705}{au}. (b)~Venus, at 1.91,
receives more than the runaway flux and even more than the empirical
recent-Venus limit (1.78): it lies inside the inner edge. Mars, at 0.43,
lies between 1.014 and 0.3438, inside the conservative habitable zone:
its dryness comes from its small size and the loss of its atmosphere, not
from its distance.
\end{solution}

\begin{solution}{exo:b3:planetary-atmospheres:4}
(a)~$\log_{10}p = 3.9895 - 443.028/93.66 = -0.741$, $p =
\qty{0.18}{bar}$. (b)~$0.05\times1.6 = \qty{0.080}{bar}$, a relative
humidity of $0.080/0.18 = 44\,\%$. (c)~The air above the seas is far from
saturated: methane evaporates from them, as water does from the Earth's
oceans, and the humid air must rise and cool before it rains.
\end{solution}

\begin{solution}{exo:b3:planetary-atmospheres:5}
$\lambda_\mathrm D = 7.02\times2.014/1.008 = 14.0$. $u_\mathrm D/u_\mathrm H
= \sqrt{1.008/2.014} = 0.71$. Ratio per atom $= 0.71\times(15.0\,
e^{-14.0})/(8.02\,e^{-7.02}) = 0.71\times\num{1.69e-3} = \num{1.2e-3}$.
\omterm{def:b3:planetary-atmospheres:jeans-escape}{Jeans escape} lets almost no deuterium out: it is a strong filter, with
$\alpha \approx 0.001$.
\end{solution}

\begin{solution}{exo:b3:planetary-atmospheres:6}
(a)~$0.016/\num{1.5575e-4} = 103$. (b)~With $\alpha = 0$ the enrichment is
$M_0/M$: $M_0 = 103\times\num{6e15} = \qty{6.2e17}{kg}$, or
$\num{6.2e17}/(1000\times\num{4.60e14}) = \qty{1.3}{m}$ of water over
Venus. (c)~$M_0/M = 103^{1/0.2} = \num{1.2e10}$, a past inventory of
\qty{7e25}{kg}, twelve times the mass of the Earth: impossible. The
enrichment needs a strongly fractionating escape (Jeans or non-thermal),
active at least during the loss of the last of the water, or a steady
state with a supply of water.
\end{solution}

\begin{solution}{exo:b3:planetary-atmospheres:7}
(a)~In four billion years ($\num{1.26e17}$~s) Mars loses
$\num{3.8e43}$ atoms, $\num{1.9e43}$ molecules of water, $\qty{5.7e17}{kg}$;
over $4\pi(\num{3.3895e6})^2 = \qty{1.44e14}{m^2}$ that is a layer
\qty{3.9}{m} deep. (b)~Twenty times more, \qty{78}{m}. (c)~Even the
high-obliquity rate reaches only the lower end of the 100 to
\qty{1500}{m}; the rest went into the crust and the ice, or was lost
faster when the Sun was young.
\end{solution}

\begin{solution}{exo:b3:planetary-atmospheres:8}
(a)~$m = pA/g = 636\times\num{1.44e14}/3.71 = \qty{2.5e16}{kg}$.
(b)~$0.1\times\num{1.26e17} = \qty{1.3e16}{kg}$, half the present
atmosphere; the whole of it in $\num{2.5e16}/0.1 = \qty{2.5e17}{s}$,
\qty{7.8}{billion~years}. (c)~The young Sun was more active and its
storms more frequent, and the loss grows during storms; other channels
(hydrogen, sputtering, photochemical escape) are not counted.
\end{solution}

\begin{solution}{exo:b3:planetary-atmospheres:9}
(a)~$S(1 - 0.193)/4 = 291$ gives $S = \qty{1442}{W.m^{-2}}$, 1.06 times
the Earth's. (b)~Venus today would absorb $0.807\times2601/4 =
\qty{525}{W.m^{-2}}$; a young Venus, $0.7\times2601 = \qty{1821}{W.m^{-2}}$
received, \qty{367}{W.m^{-2}} absorbed: both above the limit, so a
cloud-free young Venus would have run away from the start. (c)~Clouds
raise the albedo; three-dimensional models with clouds have found a path
to a habitable young Venus.
\end{solution}

\begin{solution}{exo:b3:planetary-atmospheres:10}
(a)~$H = kT/mg = \num{1.381e-23}\times1000/(16\times\num{1.661e-27}\times9.21)
= \qty{56}{km}$. (b)~$n_\mathrm{exo} = 1/(\sigma_c H) = 1/(\num{3e-19}\times
\num{5.64e4}) = \qty{5.9e13}{m^{-3}}$; $z = 200 + 56.4\ln(\num{1e16}/
\num{5.9e13}) = 200 + 56.4\times5.13 = \qty{490}{km}$. (c)~Gravity weaker
higher up makes the scale height grow with height, so the density falls
more slowly than assumed: the \omterm{def:b3:planetary-atmospheres:exobase}{exobase} moves up.
\end{solution}

\begin{solution}{exo:b3:planetary-atmospheres:11}
(a)~$\Phi = \num{3.0e13}\times\num{3e-3} = \qty{9e10}{cm^{-2}.s^{-1}} =
\qty{9e14}{m^{-2}.s^{-1}}$; over $4\pi(\num{6.371e6})^2 =
\qty{5.10e14}{m^2}$, $\num{4.6e29}$ molecules of $\ce{H2}$ per second.
(b)~An ocean holds $\num{1.4e24}/18.015\times\num{6.022e23} =
\num{4.7e46}$ molecules of water, one $\ce{H2}$ each: it lasts
$\qty{1.0e17}{s} = \qty{3.2}{billion~years}$, comparable with the age of
the Earth, which is why a mixing ratio of \num{3e-3} marks the \omterm{def:b3:planetary-atmospheres:runaway}{moist
greenhouse} limit.
\end{solution}

\begin{solution}{exo:b3:planetary-atmospheres:12}
(a)~$p_1 = g/D\kappa = 9.81/(1.66\times0.005) = \qty{1180}{Pa}$.
(b)~$1/T = 1/273.16 - \ln(1182/611.657)\times461.5/\num{2.5e6}$, $T =
\qty{282.5}{K}$, and $\sigma T^4 = \qty{361}{W.m^{-2}}$. (c)~The real limit
is lower, so the radiating level is colder and at lower pressure: real
water vapour is more opaque than this grey gas, and its effective $\kappa$
larger.
\end{solution}

\begin{solution}{pb:b3:planetary-atmospheres:1}
\textbf{1.} $g = GM/R^2 = \num{6.674e-11}\times\num{5.2e24}/(\num{6.2e6})^2 =
\qty{9.03}{m.s^{-2}}$; $\bar\rho = M/\tfrac43\pi R^3 = \qty{5210}{kg.m^{-3}}$.

\textbf{2.} $v_\mathrm{esc} = \sqrt{2GM/r} = \sqrt{2\times\num{3.47e14}/\num{6.7e6}}
= \qty{10.18}{km.s^{-1}}$.

\textbf{3.} $u = \sqrt{2kT/m}$: hydrogen \qty{4.81}{km.s^{-1}}, helium
\qty{2.41}{km.s^{-1}}, oxygen \qty{1.21}{km.s^{-1}}.

\textbf{4.} $\lambda = (v_\mathrm{esc}/u)^2$: 4.49, 17.8, 71.2.

\textbf{5.} Only hydrogen: its escaping fraction is 0.06, helium's
\num{3.5e-7}, oxygen's \num{9e-30}. None is near the hydrodynamic regime,
$\lambda \approx 1.5$ to 3.

\textbf{6.} $\lambda = 4.49\times1400/3000 = 2.09$: within the transition;
hydrogen would have streamed away hydrodynamically, perhaps dragging
heavier gases.

\textbf{7.} $(1 + 4.49)\,e^{-4.49} = 0.0618$.

\textbf{8.} $n u/2\sqrt\pi = \num{2e11}\times4806/3.545 = \qty{2.71e14}{m^{-2}.s^{-1}}$;
$\Phi_\mathrm J = \num{2.71e14}\times0.0618 = \qty{1.68e13}{m^{-2}.s^{-1}}$.

\textbf{9.} $4\pi(\num{6.7e6})^2 = \qty{5.64e14}{m^2}$; $\num{9.5e27}$ atoms per
second.

\textbf{10.} $\num{1.4e24}/18.015 = \qty{7.8e22}{mol}$, $\num{4.68e46}$ molecules,
$\num{9.36e46}$ hydrogen atoms.

\textbf{11.} $\num{9.36e46}/\num{9.5e27} = \qty{9.9e18}{s}$, about 310 billion
years.

\textbf{12.} $\lambda_\mathrm D = 4.49\times2.014/1.008 = 8.96$; escaping fraction
$9.96\,e^{-8.96} = \num{1.28e-3}$; relative efficiency $0.707\times
\num{1.28e-3}/0.0618 = 0.015$.

\textbf{13.} $H_a = kT/m_a g = \num{1.381e-23}\times200/(44.01\times\num{1.661e-27}
\times9.03) = \qty{4.18}{km}$.

\textbf{14.} $b = \qty{1.2e21}{m^{-1}.s^{-1}}$, so
\[
\Phi_\mathrm{DL} = \frac{b f(m_a - m_i)}{m_a H_a}
= \frac{\num{1.2e21}\times\num{1e-3}\times(42.0/44.0)}{4185}
= \qty{2.74e14}{m^{-2}.s^{-1}}.
\]

\textbf{15.} $\num{3.0e13}\times\num{1e-3}\times9.03/9.81 =
\qty{2.76e10}{cm^{-2}.s^{-1}} = \qty{2.76e14}{m^{-2}.s^{-1}}$: the same
within 1\,\%.

\textbf{16.} $4\pi(\num{6.2e6})^2 = \qty{4.83e14}{m^2}$; $\num{1.32e29}$ molecules
of $\ce{H2}$, $\num{2.64e29}$ atoms per second.

\textbf{17.} The diffusion limit supplies 28 times what the \omterm{def:b3:planetary-atmospheres:exobase}{exobase} lets out at
its present density. Hydrogen piles up in the thermosphere until the
Jeans flux matches the supply: in the long run the escape is
diffusion-limited, as on the Earth.

\textbf{18.} $\num{2e11}\times28 = \qty{5.6e12}{m^{-3}}$.

\textbf{19.} $\num{4.68e46}/\num{1.32e29} = \qty{3.54e17}{s}$, 11.2 billion years.

\textbf{20.} $S = 0.30/0.60^2 = 0.833$.

\textbf{21.} $d = \sqrt{L/S_\mathrm{eff}}$: inner edge $\sqrt{0.30/0.939} =
\qty{0.565}{au}$, outer edge $\sqrt{0.30/0.290} = \qty{1.02}{au}$. Kepos,
at \qty{0.60}{au}, is inside, close to the inner edge.

\textbf{22.} $0.939/0.833 = 1.13$: a brightening of 13\,\%.

\textbf{23.} $0.833\times1361\times0.7/4 = \qty{198}{W.m^{-2}}$, $T_\mathrm{eq} =
(198/\sigma)^{1/4} = \qty{243}{K}$; well below the
\qty{291}{W.m^{-2}} limit (computed for a Sun-like star, but the margin
is wide).

\textbf{24.} Three times faster: 3.7 billion years.

\textbf{25.} \textbf{11~billion~years}.
\end{solution}
